\documentclass[10pt,a4paper]{amsart}
\usepackage[margin=19mm]{geometry}
\usepackage{amsmath,amssymb,mathtools}
\usepackage[scr=rsfs,cal=euler]{mathalfa}
\usepackage{xfrac}
\usepackage[unicode,hidelinks,bookmarks=false]{hyperref}
\newcommand{\ie}{i.e.\ }
\newcommand{\cf}{cf.\ }
\newcommand{\resp}{respectively\ }
\newcommand{\Subsection}{\subsection}
\providecommand{\nobreakdash}{\nobreak\hbox{-}}
\setlength{\emergencystretch}{2em}
\multlinegap0pt
\usepackage{amssymb}
\newtheorem{lemma}{Lemma}
\newtheorem{proposition}{Proposition}
\title{On the Geometry of Complete Spacelike LW-Submanifolds in Locally Symmetric Semi-Riemannian Spaces}
\author{Jogli G. S. Ara\'ujo, Weiller F. C. Barboza}
\date{Source: arXiv:2602.14883v1 (CC BY 4.0)}
\begin{document}
\maketitle

\noindent\emph{Source and licence.} On the Geometry of Complete Spacelike LW-Submanifolds in Locally Symmetric Semi-Riemannian Spaces. Version arXiv:2602.14883v1 (CC BY 4.0), distributed by arXiv under the Creative Commons Attribution 4.0 International licence, http://creativecommons.org/licenses/by/4.0/, as recorded in the arXiv licence metadata for this exact version and shown on the abstract page https://arxiv.org/abs/2602.14883v1. Full licence notice: \texttt{licenses/arxiv-2602.14883-CC-BY-4.0.txt}.\par\smallskip

\noindent\textit{Editorial scope.} Counted unit: the article's proposition hip\_1\_prop (source0.tex lines 911--923). The article's complete original proof of it is delivered with it (source0.tex lines 925--1079, one \texttt{proof} environment of 155 lines). No mathematics is written, repaired or re-derived: the statement and the proof are the article's own lines, in the article's own notation, and no retained line is substituted or re-flowed.  Retained uncounted as the source-local context the proof needs: lines 202--204; lines 546--569; lines 610--612; lines 622--624; lines 673--683.  Nothing was omitted from any retained span, so the package carries no omission record.  No retained line contains a citation, so the wrapper carries no bibliography and no bibliography entry.  No retained line contains a figure environment, an image-inclusion command or drawing code, so no image is delivered and the package declares no asset dependency.  Editorial formatting only, in addition: an AMS \texttt{amsart} preamble that restates the article's own declarations and macros used by the retained lines, namely the environments \texttt{lemma}, \texttt{proposition} on separate counters, together with the standard packages those lines need.  The retained environments print in this document as Lemma 1, Proposition 1 in order of appearance, whereas the article prints the same items with the numbering of its own full document.  The complete native source of the article as distributed in the arXiv e-print for this exact version is bundled unchanged as \texttt{sources/194755/source0.tex}.
\begin{equation}\label{equacaoRS}
S_1 + n(n-1)R = \sum_{i,j} \overline{R}_{ijij} - n^2H^2 + S_2.
\end{equation}

\begin{lemma}\label{formula simons}
Let $X:M^{n}\looparrowright \mathbb{L}^{n+p}_{q}$ ($1\leq q\leq p$) be an submanifold with flat normal bundle and parallel normalized mean curvature vector field. Then, we get
\begin{equation}\label{simons_identity}
\begin{aligned}
\dfrac{1}{2}\Delta S \;=\;& |\nabla B|^2
  + 2\Bigg(
      \sum_{i,j,k,m,\alpha} h_{ij}^{\alpha}h_{km}^{\alpha}\overline{R}_{mijk}
      + \sum_{i,j,k,m,\alpha} h_{ij}^{\alpha}h_{jm}^{\alpha}\overline{R}_{mkik}
    \Bigg) \\[6pt]
&+ \sum_{i,j,k,\alpha,\alpha'} h_{ij}^{\alpha}h_{jk}^{\alpha'}\overline{R}_{\alpha i \alpha' k}
  + \sum_{i,j,k,\alpha,\alpha'} h_{ij}^{\alpha}h_{kk}^{\alpha'}\overline{R}_{\alpha ij \alpha'}
  + \sum_{i,j,k,\alpha,\alpha'} h_{ij}^{\alpha}h_{ij}^{\alpha'}\overline{R}_{\alpha k \alpha' k} \\[6pt]
&+ \sum_{i,j,k,\alpha,\alpha'} h_{ij}^{\alpha}h_{kj}^{\alpha'}\overline{R}_{\alpha ki \alpha'}
  + n\sum_{i,j} h_{ij}^{\,n+p-q+1}H_{ij}
  - nH\sum_{i,j,m,\alpha} h_{ij}^{\alpha}h_{mi}^{\alpha}h_{mj}^{\,n+p-q+1} \\[6pt]
&+ \sum_{\alpha,\alpha'} N(h^{\alpha}h^{\alpha'}-h^{\alpha'}h^{\alpha})
  - \sum_{\alpha,\beta} \big[\operatorname{tr}(h^{\alpha}h^{\beta})\big]^2
  - \dfrac{1}{2}\sum_{\alpha,\beta} N(h^{\alpha}h^{\beta}-h^{\beta}h^{\alpha}) \\[6pt]
&+ \sum_{\alpha,\gamma} \big[\operatorname{tr}(h^{\alpha}h^{\gamma})\big]^2
  + \dfrac{1}{2}\sum_{\alpha,\gamma} N(h^{\alpha}h^{\gamma}-h^{\gamma}h^{\alpha}) \,.
\end{aligned}
\end{equation}
where $N(A)=tr(AA^{t}),$ for all matrix $A=(a_{ij}).$
\end{lemma}

\begin{equation}\label{c1}
\overline{K}(x_i,x_{\alpha})=\dfrac{c_1}{n},\quad x_i\in \langle e_1,\cdots, e_n\rangle\quad\mbox{and}\quad x_{\alpha}\in \langle e_{n+1},\cdots,e_{n+p}\rangle,
\end{equation}

\begin{equation}\label{c3}
\overline{K}(x_{\alpha},x_{\alpha'})=\dfrac{c_3}{p},\quad x_{\alpha},x_{\alpha'}\in \langle e_{n+1},\cdots,e_{n+p}\rangle.
\end{equation}

\begin{lemma}\label{desigualdade_BeH}
Let $X:M^n\looparrowright \mathbb{L}^{n+p}_{q}$ be a complete spacelike LW-submanifold satisfying conditions \eqref{c1}-\eqref{c3} and \eqref{equacaoRS}. If the second fundamental form of $M^n$ is locally timelike and
\begin{equation}\label{H1}
(n-1)a^2+4n(\overline{\mathcal{R}}-b)\geq 0.
\end{equation}
Then, 
\begin{eqnarray}\label{desigualdade}
|\nabla B|^2 \geq n^2|\nabla H|^2.
\end{eqnarray}
Moreover, is the equality holds in ~\eqref{desigualdade} on $M^{n}$, then $H$ is constant on $M^{n}.$
\end{lemma}

\begin{proposition}
Let $X:M^n\looparrowright \mathbb{L}^{n+p}_{q}$ be a complete spacelike LW-submanifold with parallel normalized mean curvature vector field and flat normal bundle satisfying conditions \eqref{c1}-\eqref{c3}. Suppose that the second fundamental form of $M^n$ is locally timelike and there exists an orthogonal basis for $TM$ that diagonalizes simultaneously all $B_{\varepsilon},
\varepsilon\in TM^{\perp}.$ If $c=\dfrac{c_1}{n}+2c_2$ and
$(n-1)a^2+4n(\overline{\mathcal{R}}-b)\geq 0$,
then
\begin{equation}\label{hip_1_prop}
\mathcal{L}(nH)\geq |\Phi|^2\cdot P_{n,p,c,H}(|\Phi|),
\end{equation}
where 
\begin{equation}\label{hip_2_prop}
P_{n,p,c,H}(x)=-p^2x^2-\dfrac{n(n-2)}{\sqrt{n(n-1)}}Hx+n(c-H^2).
\end{equation}
\end{proposition}

\begin{proof}

Let us consider $\{e_1,e_2,\cdots,e_{n}\}$ a local orthonormal frame on $M^{n}$ such that $h_{ij}^{\alpha}=\lambda_i^{\alpha}\delta_{ij}$, for all $\alpha\in \{n+1,n+2,\cdots,n+p\}.$ From of the Lemma \ref{formula simons} we get
\begin{eqnarray*}
2\left(\sum_{i,j,k,m,\alpha}h_{ij}^{\alpha}h_{km}^{\alpha}\overline{R}_{mijk}+\sum_{i,j,k,m,\alpha}h_{ij}^{\alpha}h_{jm}^{\alpha}\overline{R}_{mkik}\right)
&=&2\sum_{i,k,\alpha}\left((\lambda_{i}^{\alpha})^2+\lambda_i^{\alpha}\lambda_k^{\alpha}\overline{R}_{kiik}\right)\\
&=&\sum_{i,k,\alpha}\overline{R}_{ikik}\left(\lambda_i^{\alpha}-\lambda_k^{\alpha}\right)^2.
\end{eqnarray*}
Since that $L_{q}^{n+p}$ satisfies the condition $\overline{K}(u,v)\geq c_2$ for any $u,v\in TM,$ we have
\begin{eqnarray*}
2\left(\sum_{i,j,k,m,\alpha}h_{ij}^{\alpha}h_{km}^{\alpha}\overline{R}_{mijk}+\sum_{i,j,k,m,\alpha}h_{ij}^{\alpha}h_{jm}^{\alpha}\overline{R}_{mkik}\right)\geq c_2\sum_{i,k,\alpha}(\lambda_i^{\alpha}-\lambda_k^{\alpha})^2=2nc_2|\Phi|^2.
\end{eqnarray*}
Now, for each $\alpha$, consider $h^{\alpha}$ the symmetric matrix $(h_{ij}^{\alpha})$, and 
\begin{eqnarray*}
S_{\alpha\beta}=\sum_{i,j}h_{ij}^{\alpha}h_{ij}^{\beta}.
\end{eqnarray*}
Then the $(p\times p)$-matrix $(S_{\alpha\beta})$ is symmetric and we can see that is diagonalizable for a choose of $e_{n+1},e_{n+2},\cdots,e_{n+p}.$ Thence 
\begin{eqnarray*}
S_{\alpha}=S_{\alpha\alpha}=\sum_{i,j}h_{ij}^{\alpha}h_{ij}^{\alpha},
\end{eqnarray*}
and we have that 
\begin{eqnarray*}
S=\sum_{\alpha}S_{\alpha}.
\end{eqnarray*}
Since that $\mathbb{L}_{q}^{n+p}$ satisfies the condition $\overline{K}(u,\eta)=\dfrac{c_1}{n},$ for any $u\in TM$ and $\eta\in TM^{\perp},$ we obtain
\begin{eqnarray*}
\sum_{i,j,k,\alpha,\alpha'}h_{ij}^{\alpha}h_{jk}^{\alpha'}\overline{R}_{\alpha i\alpha' k}&+&\sum_{i,j,k,\alpha,\alpha'}h_{ij}^{\alpha}h_{kk}^{\alpha'}\overline{R}_{\alpha ij \alpha'}+\sum_{i,j,k,\alpha,\alpha'}h_{ij}^{\alpha}h_{ij}^{\alpha'}\overline{R}_{\alpha k\alpha'k}+\sum_{i,j,k,\alpha,\alpha'}h_{ij}^{\alpha}h_{kj}^{\alpha'}\overline{R}_{\alpha ki \alpha'}\\
&=&\sum_{i,k,\alpha}(\lambda_i^{\alpha})^2\overline{R}_{\alpha k\alpha k}-nH^2c_1
=c_1|\Phi|^2.
\end{eqnarray*}




Finally note that
$$
\dfrac{1}{2}\Delta S\geq |\nabla B|^2+2nc_2|\Phi|^2+c_1|\Phi|^2+n\sum_{i,j}h_{ij}^{n+p-q+1}H_{ij}-nH\sum_{i,j,m,\alpha}h_{ij}^{\alpha}h_{mi}^{\alpha}h_{mj}^{n+p-q+1}
$$
$$
+\sum_{\alpha,\gamma}[tr(h^{\alpha}h^{\gamma})]^2-\sum_{\alpha,\beta}[tr(h^{\alpha}h^{\beta})]^2,
$$
Therefore,
\begin{equation}\label{deltaS-inequalite}
\dfrac{1}{2}\Delta S\geq |\nabla B|^2+cn|\Phi|^2+n\sum_{i,j}h_{ij}^{n+p-q+1}H_{ij}-nH\sum_{i,j,m,\alpha}h_{ij}^{\alpha}h_{mi}^{\alpha}h_{mj}^{n+p-q+1}
\end{equation}
$$
+\sum_{\alpha,\gamma}[tr(h^{\alpha}h^{\gamma})]^2-\sum_{\alpha,\beta}[tr(h^{\alpha}h^{\beta})]^2,
$$
Thence, from of $\mathcal{L}=\square+\dfrac{n-1}{2}a\Delta,$ we have
\begin{eqnarray*}
\mathcal{L}(nH)&=&\square (nH)+\dfrac{n-1}{2}a\Delta  (nH)\\
&=&\sum_{i,j}(nH\delta_{ij}-h_{ij}^{n+p-q+1})(nH)_{ij}+\dfrac{n-1}{2}a\Delta(nH)\\
&=&n^2H\sum_{i}H_{ii}-n\sum_{i,j}h_{ij}^{n+p-q+1}H_{ij}+\dfrac{n-1}{2}a\Delta (nH)\\
&=&n^2H\Delta H-n\sum_{i,j}h_{ij}^{n+p-q+1}H_{ij}+\dfrac{n-1}{2}a\Delta (nH)
\end{eqnarray*}

Note that
\begin{eqnarray*}
\Delta H^2=2H\Delta H+2|\nabla H|^2.
\end{eqnarray*}
Thus,
\begin{eqnarray*}
\mathcal{L}(nH)=\dfrac{1}{2}\Delta (n^2H^2)-n^2|\nabla H|^2-n\sum_{i,j}h_{ij}^{n+p-q+1}H_{ij}+\dfrac{n-1}{2}a\Delta(nH).
\end{eqnarray*}
Since that $R=aH+b$ and $\mathbb{L}_{q}^{n+p}$ satisfies the conditions \eqref{c1}-\eqref{c3}, we have that $\mathcal{R}$ is constant, then from \eqref{equacaoRS} we get 
\begin{eqnarray*}
\dfrac{1}{2}n(n-1)\Delta(aH)+\dfrac{1}{2}\Delta(n^2H^2)=\dfrac{1}{2}\Delta S.
\end{eqnarray*}
Therefore, using the inequality \eqref{deltaS-inequalite} and Lemma \eqref{desigualdade_BeH} we conclude that 
\begin{eqnarray}\label{retorno}
\mathcal{L}(nH)&=&\dfrac{1}{2}\Delta S-n^2|\nabla H|^2-n\sum_{i,j}h_{ij}^{n+p-q+1}H_{ij} \nonumber \\
&\geq& |\nabla B|^2-n^2|\nabla H|^2+cn|\Phi|^2-nH\sum_{i,j,m,\alpha}h_{ij}^{\alpha}h_{mi}^{\alpha}h_{mj}^{n+p-q+1}+\sum_{\alpha,\gamma}[tr(h^{\alpha}h^{\gamma})]^2-\sum_{\alpha,\beta}[tr(h^{\alpha}h^{\beta})]^2\\
&\geq& cn|\Phi|^2-nH\sum_{i,j,m,\alpha}h_{ij}^{\alpha}h_{mi}^{\alpha}h_{mj}^{n+p-q+1}+\sum_{\alpha,\gamma}[tr(h^{\alpha}h^{\gamma})]^2-\sum_{\alpha,\beta}[tr(h^{\alpha}h^{\beta})]^2. \nonumber
\end{eqnarray}
Next, we proceed with the explicit computation of the expression involving the mixed trace terms of the second fundamental forms:
\begin{equation}\label{calculo_3_termos}
-nH\sum_{\alpha}tr[h^{n+p-q+1}(h^{\alpha})^2]+\sum_{\alpha,\gamma}[tr(h^{\alpha}h^{\gamma})]^2-\sum_{\alpha,\beta}[tr(h^{\alpha}h^{\beta})]^2.
\end{equation}
Note that
\begin{eqnarray*}
nH\sum_{i,\alpha}h_{ii}^{\alpha}h_{ii}^{\alpha}h_{ii}^{n+p-q+1}&=&nH\sum_{i,\alpha}\left(\Phi_{ii}^{\alpha}+H\delta_{\alpha (n+p-q+1)}\right)^2\left(\Phi_{ii}^{n+p-q+1}+H\right)\\
&=&nH\sum_{i,\alpha}\left((\Phi_{ii}^{\alpha})^2+2H\Phi_{ii}^{\alpha}\delta_{\alpha (n+p-q+1)}+H^2\delta_{\alpha (n+p-q+1)}\right)(\Phi_{ii}^{n+p-q+1}+H)\\
&=&nH\sum_{i,\alpha}((\Phi_{ii}^{\alpha})^2\Phi_{ii}^{n+p-q+1}+H(\Phi_{ii}^{\alpha})^2+2H\Phi_{ii}^{\alpha}\Phi_{ii}^{n+p-q+1}\delta_{\alpha (n+p-q+1)}+\\
&+&2H^2\Phi_{ii}^{\alpha}\delta_{\alpha (n+p-q+1)}+H^2\Phi_{ii}^{n+p-q+1}\delta_{\alpha (n+p-q+1)}+H^3\delta_{\alpha (n+p-q+1)})\\
&=&nH\sum_{\alpha}tr[\Phi^{n+p-q+1}\cdot (\Phi^{\alpha})^2]+nH^2|\Phi|^2+n^2H^4.
\end{eqnarray*}
Moreover,

\begin{eqnarray*}
\sum_{\alpha,\beta}[tr(h^{\alpha}h^{\beta})]^2
&=&\sum_{\alpha,\beta}[\sum_{i}h_{ii}^{\alpha}h_{ii}^{\beta}]^2 
= \sum_{\alpha,\beta}[\sum_{i}(\Phi_{ii}^{\alpha}+H\delta_{\alpha (n+p-q+1)})\cdot (\Phi_{ii}^{\beta} +H\delta_{\beta (n+p-q+1)}]^2\\
&=&\sum_{\alpha,\beta}[\sum_{i}(\Phi_{ii}^{\alpha}\Phi_{ii}^{\beta}+H\Phi_{ii}^{\alpha}\delta_{\beta (n+p-q+1)}+H\Phi_{ii}^{\beta}\delta_{\alpha (n+p-q+1)}+H^2\delta_{\alpha (n+p-q+1)}\delta_{\beta (n+p-q+1)})]^2\\
&=&\sum_{\alpha,\beta}[\sum_{i}\Phi_{ii}^{\alpha}\Phi_{ii}^{\beta}+H\delta_{\beta (n+p-q+1)}\sum_{i}\Phi_{ii}^{\alpha}+H\delta_{\alpha (n+p-q+1)}\sum_{i}\Phi_{ii}^{\beta} \\
&&+H^2\delta_{\alpha (n+p-q+1)}\delta_{\beta (n+p-q+1)}\sum_{i}1]^2\\
&=&\sum_{\alpha,\beta}[tr(\Phi^{\alpha}\Phi^{\beta})+nH^2\delta_{\alpha (n+p-q+1)}\delta_{\beta (n+p-q+1)}]^2\\
&=&\sum_{\alpha,\beta}[tr(\Phi^{\alpha}\Phi^{\beta})]^2+2nH^2\sum_{\alpha,\beta}\delta_{\alpha (n+p-q+1)}\delta_{\beta (n+p-q+1)}tr(\Phi^{\alpha}\Phi^{\beta})+n^2H^4\\
&=&\sum_{\alpha,\beta}[tr(\Phi^{\alpha}\Phi^{\beta})]^2+2nH^2|\Phi^{n+p-q+1}|^2+n^2H^4,
\end{eqnarray*}
and finally
\begin{eqnarray*}
\sum_{\alpha,\gamma}[tr(h^{\alpha}h^{\gamma})]^2&=&\sum_{\alpha,\gamma}[\sum_{i}(h_{ii}^{\alpha}h_{ii}^{\gamma})]^2\\
&=&\sum_{\alpha,\gamma}[\sum_{i}(\Phi_{ii}^{\alpha}+H\delta_{\alpha (n+p-q+1)})\Phi_{ii}^{\gamma}]^2\\
&=&\sum_{\alpha,\gamma}[\sum_{i}\Phi_{ii}^{\alpha}\Phi_{ii}^{\gamma}+H\Phi_{ii}^{\gamma}\delta_{\alpha (n+p-q+1)}]^2\\
&=&\sum_{\alpha,\gamma}[tr(\Phi^{\alpha}\Phi^{\gamma})+H\delta_{\alpha (n+p-q+1)}tr(\Phi^{\gamma})]^2\\
&=&\sum_{\alpha,\gamma}[tr(\Phi^{\alpha}\Phi^{\gamma})]^2.
\end{eqnarray*}
Substituting these three equalities into \eqref{calculo_3_termos}, we obtain
\begin{align}\label{termogrande}
-nH\sum_{i,j,m,\alpha}h_{ij}^{\alpha}&h_{mi}^{\alpha}h_{mj}^{n+p-q+1}
-\sum_{\alpha,\beta}[\operatorname{tr}(h^{\alpha}h^{\beta})]^2
+\sum_{\alpha,\gamma}[\operatorname{tr}(h^{\alpha}h^{\gamma})]^2 \nonumber \\
&= -nH\sum_{\alpha}\operatorname{tr}[\Phi^{n+p-q+1}(\Phi^{\alpha})^2]
     -nH^2|\Phi|^2
     -\sum_{\alpha,\beta}[\operatorname{tr}(\Phi^{\alpha}\Phi^{\beta})]^2
     +\sum_{\alpha,\gamma}[\operatorname{tr}(\Phi^{\alpha}\Phi^{\gamma})]^2 \nonumber \\
&\quad -2nH^2|\Phi^{n+p-q+1}|^2 \nonumber \\[4pt]
&\geq -\dfrac{n(n-2)}{\sqrt{n(n-1)}}H|\Phi|^3
      +nH^2|\Phi|^2
      -2nH^2|\Phi|^2
      -\sum_{\alpha,\beta}[\operatorname{tr}(\Phi^{\alpha}\Phi^{\beta})]^2
      +\sum_{\alpha,\gamma}[\operatorname{tr}(\Phi^{\alpha}\Phi^{\gamma})]^2 \\[4pt]
&\geq -\dfrac{n(n-2)}{\sqrt{n(n-1)}}H|\Phi|^3
      -nH^2|\Phi|^2
      -\sum_{\alpha,\beta}[\operatorname{tr}(\Phi^{\alpha}\Phi^{\beta})]^2
      +\dfrac{1}{p}|\Phi|^4, \nonumber
\end{align}
furthermore, since
\begin{eqnarray*}
[tr(\Phi^{\alpha}\Phi^{\beta})]^2&=&(\Phi_{11}^{\alpha}\Phi_{11}^{\beta}+\Phi_{22}^{\alpha}\Phi_{22}^{\beta}+\cdots+\Phi_{nn}^{\alpha}\Phi_{nn}^{\beta})^2\\
&\leq&|\Phi^{\alpha}|^2\cdot |\Phi^{\beta}|^2\leq |\Phi|^2\cdot |\Phi|^2=|\Phi|^4,
\end{eqnarray*}
hence,
\begin{eqnarray*}
\sum_{\alpha,\beta}[tr(\Phi^{\alpha}\Phi^{\beta})]^2\leq |\Phi|^4\sum_{\alpha,\alpha'}1=p^2|\Phi|^4,
\end{eqnarray*}
or yet, 
\begin{eqnarray*}
-\sum_{\alpha,\beta}[tr(\Phi^{\alpha}\Phi^{\beta})]^2\geq -p^2|\Phi|^4.
\end{eqnarray*}
Therefore, from \eqref{termogrande}
\begin{eqnarray*}
-\dfrac{n(n-2)}{\sqrt{n(n-1)}}H|\Phi|^3-nH^2|\Phi|^2-\sum_{\alpha,\beta}[tr(\Phi^{\alpha}\Phi^{\beta})]^2\geq -\dfrac{n(n-2)}{\sqrt{n(n-1)}}H|\Phi|^3-nH^2|\Phi|^2-p^2|\Phi|^4+\dfrac{1}{p}|\Phi|^4.
\end{eqnarray*}
Consequently,
\begin{equation}\label{4.7}
L(nH)\geq cn|\Phi|^2-\dfrac{n(n-2)}{\sqrt{n(n-1)}}H|\Phi|^3-nH^2|\Phi|^2-p^2|\Phi|^4+\dfrac{1}{p}|\Phi|^4
\geq |\Phi|^2P_{n,p,c,H}(|\Phi|),
\end{equation}
where
\begin{equation}
P_{n,p,c,H}(x)=-p^2x^2-\dfrac{n(n-2)}{\sqrt{n(n-1)}}Hx-n(H^2-c).
\end{equation}

\end{proof}

\end{document}
